\chapter{Matriz de emprendimientos}
\label{cap:matriz}

Esta tabla reúne los fraccionamientos, clubes de campo y urbanizaciones del departamento que
este relevamiento identificó, con lo que se sabe de cada uno y con lo que no. \textbf{Los
casilleros vacíos se publican vacíos}: son la medida honesta del estado del trabajo.

La última columna consigna, cuando el acto aprobatorio se obtuvo, las condiciones que impuso
para la asignación de matrículas individuales. Es la columna que sostiene la primera tesis de
este libro, y conviene leerla en vertical.

\begin{longtable}{@{}p{2.8cm}p{1.4cm}p{3.2cm}p{5.1cm}@{}}
\toprule
\textbf{Emprendimiento} & \textbf{Munic.} & \textbf{Titularidad} & \textbf{Estado documental y exigencias} \\
\midrule
\endhead

\textbf{El Durazno} I--V y Los Duraznos 4ª \textit{(la cobertura de 2019 enumera hasta la VI; cap.\ \ref{cap:amparo})} & La Caldera & Matrícula 6.558. Héctor D'Andrea\index{D'Andrea, Héctor} y Jardín Celestial S.R.L.\ (CUIT 30-71157214-3). \textbf{Magdalena Vidal\index{Vidal, Magdalena}} figura como co-titular registral del catastro 4366 en el edicto de agua de 2015; la relación entre las matrículas 4366 y 6.558 no está establecida (cap.\ \ref{cap:ribera}) & Plano 1.164, aprobado por Res.\ 646 D/21 sólo por etapas 1, 2 y 4; el registro municipal de 2022 da por aprobadas las etapas I, III y IV y como irregulares la II y la V. \textbf{Exigencias: libre deuda y donación del espacio institucional} (8.126,79 m² = 3,05\%). En 2015 se tramitó agua para \textbf{400 lotes}, dren horizontal, carácter permanente (Expte.\ 0090034-211560/2014-0). Los arroyos El Durazno y Guaranguay tienen línea de ribera desde dic.\ 2014, \textbf{sin coordenadas publicadas}. La finca «El Durazno o Peñones» se deslindó judicialmente en 1932, a pedido de Silvano Y.\ Murúa y su esposa. Parte actora en \textit{Lazarte}; aluviones de 2018 y 2019; anegamientos en 2024 y 2026 \\
\addlinespace

\textbf{El Nogalar} I--IV & La Caldera & \textbf{Municipal}, según prensa; I y II sobre la matrícula 6557 & Consignados como loteos municipales (\textit{El Tribuno}, ene.\ 2026). Registro municipal de 2022: I y II irregulares; IV en «zona aluvional». En enero de 2026 el intendente declaró que El Nogalar 4 está dentro del lecho del río, sin catastro, y que lo desarrollaron gestiones municipales anteriores (prensa). Parte actora en \textit{Lazarte}. Clasificados como núcleos espontáneos por el art.\ 42 del COT \\
\addlinespace

\textbf{Potreros de La Caldera} & La Caldera & Matrícula 4207. Loteador convocado al taller de 2023: Ramiro Costas de Arguibel & Club de campo próximo al río. Registro municipal de 2022: irregular. Parte actora en \textit{Lazarte}. Art.\ 42 del COT \\
\addlinespace

\textbf{Solanas de La Caldera} & La Caldera & Copropiedad consignada en prensa, sin respaldo registral & Loteo cerrado. Comercializado por Cabrera Bienes Raíces; material publicado dic.\ 2024. Ingreso cedido el 4/1/2026 sobre la RN 9, km 1618. Res.\ 02/26 del Concejo pide clausura preventiva del ingreso \\
\addlinespace

\textbf{La Falda} & La Caldera & Matrícula 154; titular consignado en comunicado municipal & Km 16--17,5 RN 9, sobre ladera. Apertura a cargo de Norte Áridos. \textbf{Clausura preventiva} dic.\ 2020 por Ley 7070 y falta de autorización. El municipio determina que el sector no es urbanizable y remite al Proyecto Especial \\
\addlinespace

\textbf{Las Vertientes} & Vaqueros & Matrícula 3989, fracción 78. \textbf{Maximiliano Rodrigo Jiménez} & Club de campo (Ley 5602 y Dcto.\ 924/80), \textbf{143 ha 6.688 m²}, aprobado por Res.\ 59 D/14. \textbf{Exigencias: reglamento de condominio con servicio de agua privado, \textit{resolución de línea de ribera} y libre deuda.} La línea se determinó 67 días después, por Res.\ 80/14, \textbf{con coordenadas publicadas} \\
\addlinespace

\textbf{Chacras del Gallinato} & La Caldera & Matrícula 3.709. \textbf{El Gallinato S.R.L.} & Club de campo, \textbf{1.569 ha}, aprobado por Res.\ 585 D/15; expediente iniciado en \textbf{2007}. Desmembra 25 ha del arroyo Gallinato. \textbf{Exigencias: libre deuda, acta de recepción de EDESA, sistema de agua o consorcio, condominio de indivisión forzosa, reglamento con mínimo de 5 ha y \textit{decreto del P.E.\ de concesión de agua}} \\
\addlinespace

\textbf{Barrio La Misión} & La Caldera & Matrícula 4.065, fracción 71. \textbf{Ignacio S.\ Frías García Pinto} & Loteo de 94 ha aprobado por Res.\ 579 D/15. \textbf{Exigencia: libre deuda.} El plano \textbf{descuenta superficie por el retiro de la línea de ribera} y declara 15,4 ha de reserva no urbanizable \\
\addlinespace

\textbf{La Misión Etapa III} & La Caldera & \textbf{El Chalchal S.R.L.}, co-titular registral del catastro \textbf{3616} & Concesión de agua subterránea de pozo perforado para \textbf{580 habitantes en 116 lotes}, carácter permanente (2015). El catastro 3616 es el de la Res.\ 142/14 de línea de ribera, \textbf{sin coordenadas publicadas}, rectificada por Res.\ 236/14 \\
\addlinespace

\textbf{Finca Lesser} & Lesser \newline{\footnotesize (juris\-dicción a verificar: los actos la ubican en La Caldera, Vaqueros o San Lorenzo)} & Catastro 1869. Promotor: \textbf{Gonzalo Saravia Echevehere}, administrador judicial de la sucesión de Adolfo Saravia Valdez & «Desarrollo urbano realizado bajo la figura de \textbf{Club de Campo}». Gestión de concesión de \textbf{115 m³/día del río Lesser para abastecimiento poblacional, carácter permanente}, más riego recreativo eventual. Expediente iniciado en \textbf{2005} \\
\addlinespace

\textbf{Cerro de Buena Vista} & Vaqueros & Matrículas \textbf{1713 y 1905}. Sebastián, Nicolás, Francisco y Carlos Marcos \textbf{Gianzanti} & Urbanización con \textbf{audiencia pública ambiental} (art.\ 49 Ley 7070) el 10/10/2014 \textbf{en el Centro Deportivo Municipal de Vaqueros}, con expediente disponible en el municipio. La matrícula 1905 tiene línea de ribera del río Vaqueros \textbf{con coordenadas publicadas}. La finca fue deslindada judicialmente en \textbf{1921} a pedido de Francisco S.\ Urquiza\index{Urquiza, Francisco S.}, que en 1924 obtuvo agua del río Wierna para ella \\
\addlinespace

\textbf{Lares de la Inmaculada} & Vaqueros & Matrícula 3012. Congregación de los Hijos de la Inmaculada Concepción; fideicomiso 2017; fiduciario arq.\ Lobato & 23 ha transferidas; proyecto de 27 ha y 227 parcelas. Desarrollo a cargo de MDay. Audiencia pública abr.\ 2019 \textbf{en el SUM del colegio del proponente}. Amparo rechazado nov.\ 2019 \\
\addlinespace

\textbf{Loteo con sanción vigente} & La Caldera &  & Referido por el centro vecinal de La Calderilla, ene.\ 2026: habría desmontado monte ribereño sin autorización y sido multado por Ambiente, y el municipio habría avanzado igual con audiencia pública. \textbf{No verificado}: el buscador de avisos del Boletín Oficial no devuelve ninguna convocatoria a audiencia pública en el municipio entre 2019 y 2026 (cap.\ \ref{cap:infraestructura}) \\
\addlinespace

\textbf{Ayres de Quitilipi} & La Caldera & Matrícula 5831. Loteador convocado al taller de 2023: Augusto Torino & Plano 1078. Intimado por el municipio en 2020, presentó en noviembre \textbf{certificado de no inundabilidad} (expte.\ 0090034-248087/2014) y \textbf{certificado de aptitud ambiental} provincial (expte.\ 0090227-310880/2018): \textbf{es el único de esta matriz con constancia de ambos} \\
\addlinespace

\textbf{La Milagrosa} & La Caldera & Matrícula 5800. Loteador convocado al taller: Matías Marcelo Cornejo Isasmendi & Registro municipal de 2022: irregular \\
\addlinespace

\textbf{Valle Alegre} & La Caldera & Loteador convocado al taller: Luis Cabanilla & Consorcio de agua pública (expte.\ 0090034-86490), con asambleas en 2024 y 2026. Intimado por el municipio en 2020 \\
\addlinespace

\textbf{Vistas del Lago} & La Caldera & Loteador convocado al taller: Federico Antonio Fernández & Intimado por el municipio en 2020 a presentar planos y certificados \\
\addlinespace

\textbf{Santiago Apóstol} & La Caldera &  & Desbordes en 2024 y 2026; defensas en ejecución en enero de 2026; ampliación de la red de agua anunciada en septiembre de 2026 \\
\addlinespace

\textbf{Loteo de las matrículas 1555 y 161} & La Caldera &  & Loteo abierto con sumario ambiental por desmonte ilegal (expte.\ 0090227-229822/2022-0); denuncia penal por incendio en el catastro 1555 (2022) \\
\addlinespace

\textbf{Los Maitines} & San Lorenzo & Desarrollador consignado en prensa & \textbf{Caso comparado.} Zona amarilla Cat.\ II (Ley 26.331). Peritaje judicial nov.\ 2023 \\
\bottomrule
\end{longtable}

\section*{Lo que la última columna muestra}

Ordenados por fecha del acto aprobatorio, los cuatro fraccionamientos cuyo texto se obtuvo
recibieron condiciones muy distintas para la asignación de matrículas.

{\small
\begin{longtable}{@{}p{3.6cm}p{1.3cm}p{7.6cm}@{}}
\toprule
\textbf{Emprendimiento} & \textbf{Año} & \textbf{Condiciones} \\
\midrule
\endhead
Las Vertientes & 2014 & Agua privada, \textbf{línea de ribera}, libre deuda \\
Barrio La Misión & 2015 & Libre deuda \textit{(el plano ya descuenta la ribera)} \\
Chacras del Gallinato & 2015 & Seis condiciones, incluido decreto de agua \\
\textbf{El Durazno} & \textbf{2021} & \textbf{Libre deuda y donación institucional} \\
\bottomrule
\end{longtable}}

El fraccionamiento cuyos arroyos se desbordaron casi tres años antes de su aprobación es el que
menos condiciones recibió. \textbf{Los regímenes son distintos y de distinto rango}: los clubes
de campo se rigen por la \textbf{Ley 5602/80 y su Decreto 924/80}; las urbanizaciones privadas,
por una \textbf{resolución de la Junta de Catastro}; y los loteos abiertos, como El Durazno, por el \textbf{Título V de la Ley 2.308} y su reglamentación de 2019, que exige, entre otros, certificados de no inundabilidad y de aptitud ambiental. El capítulo \ref{cap:loteo} desarrolla qué explica esa
diferencia y qué no.
Pero el Barrio La Misión es, como El Durazno, un loteo abierto, aprobado seis años antes ---bajo el Decreto 1.410/73 y por el entonces Ministerio de Economía---, y su plano descuenta la línea de ribera.

\section*{Campos a completar, según los artículos del COT de La Caldera}

Para cada emprendimiento: dictamen del Órgano Técnico de Aplicación (art.\ 22); Certificado
de Uso Conforme (art.\ 137); Certificado de Prefactibilidad (art.\ 138); Certificado de
Aptitud Ambiental Municipal (art.\ 140); Certificado de Aptitud Técnica (art.\ 141); Permiso
de Inicio de Obra de Urbanización (art.\ 143); Garantía de Ejecución de Obras (art.\ 144);
Licencia de Comercialización (art.\ 145); actas de Control de Avance (art.\ 146); Certificado
de Final de Obra (art.\ 147).

\section*{Campos registrales}

Matrícula catastral; superficie de la parcela madre; \textbf{cadena de titularidad}, que para
las fincas históricas puede remontarse a los deslindes judiciales de 1921 y 1932 y al catastro
provincial de 1929; instrumento de estructuración; número de expediente ante la Secretaría de
Ambiente; categoría OTBN del predio; fecha y número de la ordenanza municipal de aprobación;
votación nominal en el Concejo; obras de infraestructura comprometidas y su estado; cantidad
de parcelas escrituradas.
